\medskip

Soit $g$ la fonction définie sur l'intervalle $]0~;~+\infty[$ par 
$g(x) = 1+ x^2[1 - 2 \ln (x)].$

La fonction $g$ est dérivable sur l'intervalle $]0~;~+\infty[$ et on note $g'$ sa fonction dérivée.

On appelle $\mathcal{C}$ la courbe représentative de la fonction $g$ dans un repère orthonormé du plan.

\medskip

\textbf{PARTIE A}



\begin{enumerate}
\item $g(\text{e}) = 1+\e^2\left (1-2\ln(\e)\right ) = 1+\e^2\left (1-2\right ) 
=1-\e^2$

$\e>2$ donc $\e^2>4$ donc $1-\e^2<0$, et donc $g(\e)<0$.

\item 
$\ds\lim_{x \to +\infty} \ln(x)=+\infty$ donc $\ds\lim_{x \to +\infty} 1-2\ln(x) = -\infty$ donc $\ds\lim_{x \to +\infty} x^2\left (1-2\ln(x)\right )=-\infty$.

On en déduit que $\ds\lim_{x \to + \infty} g(x) = - \infty$.

\item 
	\begin{enumerate}
		\item 
Pour $x\in\;]0\;; +\infty[$, on a:	\\
$g'(x)=2x\left (1-2\ln(x)\right ) + x^2\left (-\dfrac{2}{x}\right )
= 2x -4x\ln(x) -2x=-4x\ln(x)$.	
		
		\item Pour étudier le sens de variation de la fonction $g$ sur l'intervalle $]0~;~+\infty[$, on détermine le signe de $g'(x)$ sur cet intervalle.
		
\begin{center}
{\renewcommand{\arraystretch}{1.5}
\def\esp{\hspace*{2cm}}
$\begin{array}{|c | l*{4}{c} |} 
\hline
x  & 0 & \esp & 1 & \esp  & +\infty \\
\hline
-4x &  & \pmb{-} &  \vline\hspace{-2.7pt}{\phantom 0} & \pmb{-} &   \\
\hline
\ln(x) &\vline\:\vline  & \pmb{-} &  \vline\hspace{-2.7pt}{0} & \pmb{+} &   \\
\hline
g'(x)=-4x\ln(x) & \vline\:\vline & \pmb{+} &  \vline\hspace{-2.7pt}{0} & \pmb{-} &   \\
\hline
\end{array}$
}
\end{center}
		
\begin{list}{\textbullet}{Donc la fonction $g$}
\item est strictement croissante sur $]0\;; 1[$;
\item est strictement décroissante sur $]1\;; +\infty[$;
\item admet en $x=1$  un maximum égal à $g(1)=1+1^2\left (1-\ln(1)\right )=2$.
\end{list}		

		\item
On trace le tableau des variations de la fonction $g$ sur $[1~;~+\infty[$.

\begin{center}
{\renewcommand{\arraystretch}{1.3}
\psset{nodesep=3pt,arrowsize=2pt 3}
\def\esp{\hspace*{2.5cm}}
\def\hauteur{0pt}
$\begin{array}{|c|l*4{c}|}
\hline
x & 0  & \esp & 1 & \esp & +\infty \\ 
\hline
g'(x) & \vline\;\vline &   \pmb{+} & \vline\hspace{-2.7pt}0 & \pmb{-} & \\ 
\hline
 &  \vline\;\vline & &   \Rnode{max}{2}  &  &   \\  
g(x) &\vline\;\vline &     &  &  &  \rule{0pt}{\hauteur} \\ 
 &\vline\;\vline \Rnode{min1}{\phantom{0}} &   &  &  &   \Rnode{min2}{-\infty} \rule{0pt}{\hauteur}    
 \ncline{->}{min1}{max} 
 \ncline{->}{max}{min2} 
 \\ 
\hline
\end{array} $
}
\end{center}	

D'après ce tableau de variations, on peut dire que l'équation $g(x)=0$ admet une solution unique sur l'intervalle $]1\;; +\infty[$. On appelle $\alpha$ cette solution.

		\item 
$\left .
\begin{array}{l}
g(1) =2>0\\
g(2)\approx -0,55<0
\end{array}		
\right \}$ donc $\alpha \in\;[1\;; 2]$

$\left .
\begin{array}{l !{\approx} l}
g(1,8) & 0,43>0\\
g(1.9) &  -0,024<0
\end{array}		
\right \}$ donc $\alpha \in\;[1,8\;; 1,9]$
		
$\left .
\begin{array}{l !{\approx} l}
g(1,89) & 0,024>0\\
g(1.90) &  -0,024<0
\end{array}		
\right \}$ donc $\alpha \in\;[1,89\;; 1,90]$			
	\end{enumerate}

\item On en déduit le signe de la fonction $g$ sur l'intervalle $[1~;~+\infty[$.

\begin{center}
\renewcommand{\arraystretch}{1.3}
\def\esp{\hspace*{2cm}}
$\begin{array}{|c | *{5}{c} |} 
\hline
x  & 1 & \esp & \alpha & \esp  & +\infty \\
\hline
g(x) &  & \pmb{+} &  \vline\hspace{-2.7pt}{0} & \pmb{-} &    \\
\hline
\end{array}$
\renewcommand{\arraystretch}{1}
\end{center}

\end{enumerate}

\medskip

\textbf{PARTIE B}

\medskip

\begin{enumerate}
\item On admet que, pour tout $x$ appartenant à l'intervalle $[1~;~\alpha]$, $g''(x) = - 4[\ln (x) + 1]$.



Sur $[1~;~\alpha]$, $\ln(x)\geqslant 0$ donc $\ln(x)+1>0$ donc $-4\left (\ln(x)+1\right )<0$. \\
On en déduit que $g''(x)>0$ et donc que la fonction $g$ est concave sur $[1~;~\alpha]$.
\end{enumerate}


\begin{enumerate}[resume]
\item 
\begin{multicols}{2}
Sur la figure ci-contre, A et B sont les points de la courbe $\mathcal{C}$ d'abscisses respectives 1 et $\alpha$.

\begin{center}
\psset{unit=1cm,arrowsize=2pt 3}
\begin{pspicture}(-0.2,-0.2)(3,3)
\psgrid[gridlabels=0pt,subgriddiv=1,gridwidth=0.15pt]
\psaxes[linewidth=1.25pt,labelFontSize=\scriptstyle]{->}(0,0)(0,0)(3,3)
\psplot[plotpoints=1000,linewidth=1.25pt,linecolor=red]{1}{1.9}{1 x ln 2 mul sub x dup mul mul 1 add}
\uput[ul](1,2){\red A}\uput[ur](1.9,0){\red B}\uput[d](1.9,0.05){$\alpha$}
\uput[ur](1.5,1.5){\red $\mathcal{C}$}
\end{pspicture}
\end{center}

\end{multicols}

	\begin{enumerate}
		\item La droite (AB) a pour équation réduite:
		
$\dfrac{y-y_{\text A}}{x - x_{\text A}} = \dfrac{y_{\text B}-y_{\text A}}{x_{\text B} - x_{\text A}}
\iff
\dfrac{y-2}{x-1} = \dfrac{0-2}{\alpha - 1}
\iff
y=\dfrac{-2 \left (x-1\right )}{\alpha - 1} +2$

$\iff
y= \dfrac{-2}{\alpha - 1} x +\dfrac{2}{\alpha - 1} + \dfrac{2\alpha - 2}{\alpha - 1}
\iff
y= \dfrac{-2}{\alpha - 1} x +\dfrac{2\alpha}{\alpha - 1}$
		
		\item 
Sur l'intervalle $[1~;~\alpha]$, la fonction est concave, donc sa courbe représentative est située au dessus de toute sécante, donc au dessus du segment [AB].

On en déduit que sur $[1~;~\alpha]$, on a: $g(x) \geqslant \dfrac{- 2}{\alpha - 1} x + \dfrac{2\alpha}{\alpha - 1}$.
		
		
	\end{enumerate}
\end{enumerate}
